% =====================================================================
%  PREÁMBULO — Curso de Ingeniería FUSAL
%  Todo lo que configura el documento vive aquí. main.tex se queda limpio.
% =====================================================================

% ---- Idioma y codificación -------------------------------------------
\usepackage[utf8]{inputenc}
\usepackage[T1]{fontenc}
\usepackage[spanish,es-noshorthands]{babel}   % es-noshorthands evita choques con listings/tikz/siunitx
\usepackage{csquotes}

% ---- Tipografía ------------------------------------------------------
\usepackage{amsmath,amssymb,mathtools}
\usepackage{newtxtext}          % Texto tipo Times
\usepackage{newtxmath}          % Matemáticas a juego
\usepackage{microtype}
% Para cambiar de fuente: sustituye las dos líneas newtx por
%   \usepackage{XCharter}\usepackage[charter]{mathdesign}   (Charter, más legible en pantalla)
%   \usepackage{palatino}                                    (Palatino)

% ---- Geometría de página --------------------------------------------
\usepackage[a4paper,left=3cm,right=2.5cm,top=3cm,bottom=3cm,headheight=15pt]{geometry}
\usepackage{setspace}
\onehalfspacing
\setlength{\parskip}{0.4em}

% ---- Colores del equipo ---------------------------------------------
\usepackage{xcolor}
\definecolor{fusalPrim}{HTML}{1F3864}   % azul oscuro  <- cambia aquí el color corporativo
\definecolor{fusalSec}{HTML}{C00000}    % rojo acento
\definecolor{fusalGris}{HTML}{6E6E6E}
\definecolor{fusalClaro}{HTML}{EEF2F8}

% ---- Figuras y tablas ------------------------------------------------
\usepackage{graphicx}
\graphicspath{{Imagenes/}{Imagenes/dinamica/}{Imagenes/suspension/}%
              {Imagenes/chasis/}{Imagenes/aero/}{Imagenes/refrigeracion/}%
              {Imagenes/transmision/}{Imagenes/transversal/}}
\usepackage{float}
\usepackage{subcaption}
\usepackage[font=small,labelfont=bf,labelsep=period]{caption}
\usepackage{booktabs,tabularx,longtable,multirow,array}
\usepackage{rotating}
\usepackage{pdfpages}
\usepackage{enumitem}
\setlist{itemsep=2pt,topsep=4pt}

% ---- Unidades --------------------------------------------------------
\usepackage{siunitx}
\sisetup{
  output-decimal-marker = {,},   % coma decimal (convenio español)
  per-mode              = symbol,
  separate-uncertainty  = true,
  range-phrase          = {--},
  range-units           = single,
  detect-weight         = true
}
\DeclareSIUnit{\rpm}{rpm}
\DeclareSIUnit{\bar}{bar}
\DeclareSIUnit{\gacc}{g}

% ---- Encabezados y pies ---------------------------------------------
\usepackage{fancyhdr}
\pagestyle{fancy}
\fancyhf{}
\fancyhead[L]{\small\textcolor{fusalGris}{\nouppercase{\leftmark}}}
\fancyhead[R]{\small\textcolor{fusalGris}{\thepage}}
\renewcommand{\headrulewidth}{0.4pt}
\fancypagestyle{plain}{\fancyhf{}\fancyfoot[C]{\small\thepage}\renewcommand{\headrulewidth}{0pt}}

% ---- Estilo de títulos ----------------------------------------------
\usepackage{titlesec}
\titleformat{\chapter}[hang]
  {\normalfont\sffamily\huge\bfseries\color{fusalPrim}}
  {\thechapter\hspace{10pt}\textcolor{fusalGris}{\textbar}\hspace{10pt}}{0pt}{\huge}
\titlespacing*{\chapter}{0pt}{-15pt}{28pt}
\titleformat{\section}
  {\normalfont\sffamily\Large\bfseries\color{fusalPrim}}{\thesection}{1em}{}
\titleformat{\subsection}
  {\normalfont\sffamily\large\bfseries\color{fusalGris}}{\thesubsection}{1em}{}

% ---- Índices ---------------------------------------------------------
\usepackage{makeidx}
\makeindex
\setcounter{tocdepth}{1}      % en el índice general, solo hasta \section
\setcounter{secnumdepth}{3}

% ---- Código fuente (carpeta Codigo/) ---------------------------------
\usepackage{listings}
\definecolor{codekw}{HTML}{0033B3}
\definecolor{codecmt}{HTML}{5C8A4A}
\definecolor{codestr}{HTML}{A31515}
\definecolor{codebg}{HTML}{F7F7F7}
\definecolor{codenum}{HTML}{999999}
\lstdefinestyle{cursocode}{
    basicstyle=\ttfamily\scriptsize,
    keywordstyle=\color{codekw}\bfseries,
    commentstyle=\color{codecmt}\itshape,
    stringstyle=\color{codestr},
    numbers=left, numberstyle=\tiny\color{codenum}, numbersep=7pt,
    backgroundcolor=\color{codebg},
    frame=single, rulecolor=\color{codenum!60},
    breaklines=true, breakatwhitespace=false, showstringspaces=false,
    tabsize=4, columns=fullflexible, keepspaces=true,
    extendedchars=true, inputencoding=utf8, literate=%
        {á}{{\'a}}1 {é}{{\'e}}1 {í}{{\'i}}1 {ó}{{\'o}}1 {ú}{{\'u}}1
        {ñ}{{\~n}}1 {Á}{{\'A}}1 {É}{{\'E}}1 {Í}{{\'I}}1 {Ó}{{\'O}}1
        {Ú}{{\'U}}1 {Ñ}{{\~N}}1 {¿}{{?`}}1 {¡}{{!`}}1 {°}{{\textdegree}}1
        {²}{{\textsuperscript{2}}}1 {³}{{\textsuperscript{3}}}1
        {µ}{{\textmu}}1 {·}{{\textperiodcentered}}1 {×}{{\texttimes}}1
        {–}{{--}}1 {—}{{---}}1 {•}{{\textbullet}}1
}
\lstset{style=cursocode,inputpath=Codigo/}
\renewcommand{\lstlistingname}{Código}
\renewcommand{\lstlistlistingname}{Índice de códigos}

% ---- Dibujo ----------------------------------------------------------
\usepackage{tikz}
\usetikzlibrary{arrows.meta,calc,positioning,patterns,decorations.pathmorphing,decorations.pathreplacing}
\usepackage{pgfplots}
\usepgfplotslibrary{groupplots,fillbetween}
\pgfplotsset{compat=1.18}

% Estilo común de todas las gráficas del libro. Úsalo siempre:
%   \begin{axis}[cursoplot, xlabel=..., ylabel=...]   una gráfica a ancho de página
%   \begin{axis}[cursopanel, ...]                     media página (dos en paralelo)
\pgfplotsset{/pgf/number format/.cd, use comma, 1000 sep={}}
\pgfplotsset{
  cursoplot/.style={
    width=0.92\linewidth, height=6.2cm,
    scale only axis, axis lines=left,
    label style={font=\sffamily\small},
    tick label style={font=\sffamily\scriptsize},
    title style={font=\sffamily\small\bfseries,color=fusalPrim},
    legend style={font=\sffamily\scriptsize,draw=fusalGris!50,fill=white,
                  fill opacity=0.88,text opacity=1,rounded corners=1pt},
    legend cell align=left,
    grid=both, major grid style={draw=fusalGris!25},
    minor grid style={draw=fusalGris!12},
    scaled ticks=false, tick align=outside,
    every axis plot/.append style={line width=0.9pt},
    clip mode=individual},
  % Panel para dos gráficas en paralelo. Aquí «width» es el ancho TOTAL del
  % dibujo (rótulos incluidos), por eso lleva scale only axis=false.
  cursopanel/.style={cursoplot, scale only axis=false,
    width=0.47\linewidth, height=6.0cm},
  cursotercio/.style={cursoplot, scale only axis=false,
    width=0.32\linewidth, height=5.4cm,
    label style={font=\sffamily\scriptsize},
    tick label style={font=\sffamily\tiny},
    title style={font=\sffamily\scriptsize\bfseries,color=fusalPrim},
    legend style={font=\sffamily\tiny,draw=fusalGris!50,fill=white,
                  fill opacity=0.88,text opacity=1,rounded corners=1pt}},
}
% Magic Formula de Pacejka, para dibujar curvas de neumático (capítulos 8 y 9):
%   mf(x,B,C,D,E) = D sin( C arctan( (1-E) Bx + E arctan(Bx) ) )
% OJO con los grados: en pgfmath atan() devuelve grados y sin() los espera, por
% eso la arcotangente interior lleva rad(). Va aquí y no en el capítulo porque
% declare function es global y no se puede declarar dos veces.
\tikzset{
  declare function={
    mf(\x,\B,\C,\D,\E) = \D*sin(\C*atan((1-\E)*\B*\x + \E*rad(atan(\B*\x))));
  }
}

% Colores de curva. El «cycle list» automático de pgfplots no es fiable en todas
% las instalaciones, así que cada \addplot lleva su estilo explícito:
%   \addplot[curvaA] ... \addplot[curvaB] ...
\tikzset{
  curvaA/.style={fusalPrim,thick},
  curvaB/.style={fusalSec,thick},
  curvaC/.style={teal!65!black,thick},
  curvaD/.style={orange!85!black,thick},
  curvaE/.style={violet!65!black,thick},
  curvaAux/.style={fusalGris,thick,densely dashed},
  cursonota/.style={font=\sffamily\scriptsize,align=center,text=fusalGris},
  cursomarca/.style={fusalSec,densely dashed,line width=0.6pt},
}

% ---- Bibliografía ----------------------------------------------------
\usepackage[style=numeric-comp,backend=biber,sorting=none,maxbibnames=99]{biblatex}
\addbibresource{bibliografia.bib}

% ---- Notas de trabajo (modo borrador) --------------------------------
% Con \borradortrue (en main.tex) se muestran los \todo y los avisos.
\newif\ifborrador
\usepackage[colorinlistoftodos,prependcaption,textsize=footnotesize,spanish]{todonotes}
\setlength{\marginparwidth}{2cm}   % silencia el aviso de todonotes (solo usamos todos en línea)
\newcommand{\pendiente}[1]{\ifborrador\todo[inline,color=yellow!40]{#1}\fi}
\newcommand{\falta}[1]{\ifborrador\todo[inline,color=red!25]{FALTA: #1}\fi}
\newcommand{\figpendiente}[1]{\ifborrador\missingfigure{#1}\fi}

% ---- Hiperenlaces (siempre lo último salvo cleveref) -----------------
\usepackage[hidelinks,bookmarksnumbered]{hyperref}
\hypersetup{
  colorlinks = true,
  linkcolor  = fusalPrim,
  citecolor  = fusalSec,
  urlcolor   = fusalPrim,
  pdftitle   = {Curso de Ingeniería FUSAL — Mechanical},
  pdfauthor  = {Formula Student USAL}
}
\usepackage[nameinlink]{cleveref}
\crefname{figure}{figura}{figuras}
\crefname{table}{tabla}{tablas}
\crefname{equation}{ecuación}{ecuaciones}
\crefname{chapter}{capítulo}{capítulos}
\crefname{section}{apartado}{apartados}

% ---- Entornos propios y macros --------------------------------------
% =====================================================================
%  ENTORNOS PROPIOS
%  Los seis cuadros que estructuran cada capítulo. Úsalos siempre igual:
%  el lector debe poder hojear el libro buscando solo los recuadros.
% =====================================================================
\usepackage[most]{tcolorbox}
\tcbuselibrary{breakable,skins}

% --- Plantilla base ---------------------------------------------------
\tcbset{
  cursobase/.style={
    breakable, enhanced, sharp corners=downhill,
    boxrule=0pt, leftrule=3pt, toprule=0pt, bottomrule=0pt, rightrule=0pt,
    left=8pt, right=8pt, top=6pt, bottom=6pt,
    fonttitle=\sffamily\bfseries\small,
    coltitle=black, attach title to upper=\par\medskip
  }
}

% --- 1. Idea clave ----------------------------------------------------
\newtcolorbox{clave}[1][]{cursobase,
  colback=fusalClaro, colframe=fusalPrim,
  title={IDEA CLAVE}, #1}

% --- 2. Caso real del equipo -----------------------------------------
\newtcolorbox{caso}[1][]{cursobase,
  colback=fusalPrim!5, colframe=fusalPrim!70,
  title={CASO FUSAL\ifx\relax#1\relax\else\ --- #1\fi}}

% --- 3. Referencia al reglamento -------------------------------------
\newtcolorbox{regla}[1][]{cursobase,
  colback=gray!8, colframe=fusalGris,
  title={DEL REGLAMENTO\ifx\relax#1\relax\else\ --- #1\fi}}

% --- 4. Errores frecuentes -------------------------------------------
\newtcolorbox{ojo}[1][]{cursobase,
  colback=fusalSec!5, colframe=fusalSec,
  title={OJO --- ERRORES FRECUENTES}, #1}

% --- 5. Entregable de la sesión --------------------------------------
\newtcolorbox{entregable}[1][]{cursobase,
  colback=green!5, colframe=green!45!black,
  title={ENTREGABLE}, #1}

% --- 6. Preguntas de los jueces --------------------------------------
\newtcolorbox{juez}[1][]{cursobase,
  colback=orange!6, colframe=orange!70!black,
  title={LO QUE TE VA A PREGUNTAR UN JUEZ}, #1}

% --- Objetivos de aprendizaje (cabecera de cada capítulo) -------------
\newenvironment{objetivos}
  {\begin{tcolorbox}[cursobase, colback=white, colframe=fusalPrim,
      title={AL TERMINAR ESTE CAPÍTULO DEBES SABER}]
   \begin{itemize}[leftmargin=*,label=\textcolor{fusalPrim}{\textbullet}]}
  {\end{itemize}\end{tcolorbox}}

% --- Ficha de sesión (para impartir el capítulo como clase) -----------
% Uso: \fichasesion{Duración}{Requisitos previos}{Material necesario}
\newcommand{\fichasesion}[3]{%
  \begin{tcolorbox}[cursobase, colback=gray!5, colframe=fusalGris,
      title={FICHA DE SESIÓN}]
  \begin{description}[leftmargin=!,labelwidth=2.6cm,font=\sffamily\bfseries\small]
    \item[Duración] #1
    \item[Requisitos] #2
    \item[Material] #3
  \end{description}
  \end{tcolorbox}}

% =====================================================================
%  MACROS
%  Símbolos que se repiten en todo el libro. Escribe SIEMPRE con macro,
%  nunca a mano: si un día cambias el convenio, cambias solo esta hoja.
% =====================================================================

% --- Generales --------------------------------------------------------
\usepackage{xspace}
\newcommand{\FUSAL}{\textsc{FUSAL}\xspace}
\newcommand{\etal}{\textit{et al.}\xspace}
\newcommand{\dif}{\mathrm{d}}
\newcommand{\deriv}[2]{\frac{\dif #1}{\dif #2}}
\newcommand{\pderiv}[2]{\frac{\partial #1}{\partial #2}}
\newcommand{\vect}[1]{\mathbf{#1}}
\newcommand{\prom}[1]{\overline{#1}}
% Texto guía dentro de los recuadros: qué hay que escribir ahí. Bórralo al redactar.
\newcommand{\redactar}[1]{\textcolor{fusalGris}{\itshape #1}}

% --- Dinámica vehicular ----------------------------------------------
\newcommand{\Fz}{F_z}                       % carga vertical
\newcommand{\Fy}{F_y}                       % fuerza lateral
\newcommand{\Fx}{F_x}                       % fuerza longitudinal
\newcommand{\Mz}{M_z}                       % par autoalineante
\newcommand{\slipa}{\alpha}                 % ángulo de deriva
\newcommand{\slipr}{\kappa}                 % deslizamiento longitudinal
\newcommand{\camber}{\gamma}                % caída
\newcommand{\muy}{\mu_y}
\newcommand{\mux}{\mu_x}
\newcommand{\ay}{a_y}
\newcommand{\ax}{a_x}
\newcommand{\CdG}{\mathrm{CdG}}             % centro de gravedad
\newcommand{\hcg}{h_{\CdG}}
\newcommand{\LLTD}{\mathrm{LLTD}}
\newcommand{\Kus}{K_{us}}                   % gradiente de subviraje
\newcommand{\batalla}{L}
\newcommand{\via}{t}
\newcommand{\Kroll}{K_{\phi}}               % rigidez a balanceo

% --- Neumático y su modelado (capítulos 8 y 9) ------------------------
\newcommand{\Calpha}{C_{\slipa}}            % rigidez de deriva
\newcommand{\Ckappa}{C_{\slipr}}            % rigidez de deslizamiento
\newcommand{\tneu}{t_p}                     % avance neumático
\newcommand{\pneu}{p_i}                     % presión de inflado
\newcommand{\Tneu}{T_n}                     % temperatura del neumático
\newcommand{\Fzo}{F_{z0}}                   % carga vertical de referencia
\newcommand{\dfz}{\mathrm{d}f_z}            % variación normalizada de carga
\newcommand{\lamu}{\lambda_{\mu}}           % factor de escala de agarre
% Coeficientes de la Magic Formula. OJO: \Cp ya es el coeficiente de
% presión aerodinámico, por eso estos llevan prefijo MF.
\newcommand{\MFB}{B}                        % factor de rigidez
\newcommand{\MFC}{C}                        % factor de forma
\newcommand{\MFD}{D}                        % valor de pico
\newcommand{\MFE}{E}                        % factor de curvatura
\newcommand{\MFSh}{S_h}                     % desplazamiento horizontal
\newcommand{\MFSv}{S_v}                     % desplazamiento vertical

% --- Suspensión -------------------------------------------------------
\newcommand{\MR}{\mathrm{MR}}               % relación de instalación
\newcommand{\Kw}{K_w}                       % rigidez en rueda
\newcommand{\Ks}{K_s}                       % rigidez del muelle
\newcommand{\fr}{f_r}                       % frecuencia de suspensión

% --- Aerodinámica -----------------------------------------------------
\newcommand{\Cl}{C_L}
\newcommand{\Cd}{C_D}
\newcommand{\Cp}{C_p}
\newcommand{\ClA}{C_L A}
\newcommand{\CdA}{C_D A}
\newcommand{\qdin}{q_{\infty}}              % presión dinámica
\newcommand{\Vinf}{V_{\infty}}
\newcommand{\Rey}{\mathit{Re}}
\newcommand{\CoP}{\mathrm{CoP}}             % centro de presiones

% --- Transferencia de calor ------------------------------------------
\newcommand{\UA}{UA}
\newcommand{\NTU}{\mathit{NTU}}
\newcommand{\efec}{\varepsilon}             % eficacia del intercambiador
\newcommand{\Nus}{\mathit{Nu}}
\renewcommand{\Pr}{\mathit{Pr}}   % redefine el \Pr de LaTeX (probabilidad), aquí es Prandtl
\newcommand{\Stan}{\mathit{St}}
\newcommand{\jcol}{j}                       % factor de Colburn
\newcommand{\ffan}{f}                       % factor de fricción
\newcommand{\mdot}{\dot{m}}
\newcommand{\Qdot}{\dot{Q}}
\newcommand{\Dh}{D_h}                       % diámetro hidráulico
\newcommand{\LMTD}{\Delta T_{\mathrm{ml}}}

% --- Transmisión ------------------------------------------------------
\newcommand{\itot}{i_{\mathrm{tot}}}        % relación total
\newcommand{\Tmot}{T_{\mathrm{mot}}}
\newcommand{\rdin}{r_{\mathrm{din}}}        % radio dinámico

% =====================================================================
%  ESTILOS DE DIAGRAMA DE FLUJO
%  Usados en los organigramas del libro. Convenio:
%    proceso    rectángulo      una acción
%    decision   rombo           una pregunta con salidas Sí/No
%    terminal   elipse          principio o final del flujo
%    nota       recuadro gris   detalle de lo que incluye un paso
%    conector   triángulo       enlace a un subproceso desarrollado aparte
%    propuesto  trazo discontinuo rojo: añadido, todavía no acordado
% =====================================================================
\usetikzlibrary{shapes.geometric}

\tikzset{
  proceso/.style={
    draw=fusalPrim!75, fill=fusalPrim!6, rounded corners=2pt,
    align=center, inner sep=4pt, minimum height=8mm, text width=26mm,
    font=\sffamily\scriptsize},
  decision/.style={
    draw=fusalPrim!75, fill=white, diamond, aspect=1.7,
    align=center, inner sep=0pt, text width=19mm,
    font=\sffamily\scriptsize},
  terminal/.style={
    draw=fusalPrim, fill=fusalPrim!15, ellipse,
    align=center, inner sep=3pt, text width=28mm,
    font=\sffamily\scriptsize\bfseries},
  nota/.style={
    draw=fusalGris, fill=gray!7, align=left, inner sep=4pt,
    text width=32mm, font=\sffamily\tiny},
  conector/.style={
    draw=fusalSec, fill=fusalSec!12, regular polygon,
    regular polygon sides=3, inner sep=0pt, minimum size=7mm,
    font=\sffamily\tiny\bfseries, text=fusalSec},
  propuesto/.style={
    draw=fusalSec, fill=fusalSec!5, dashed, dash pattern=on 2pt off 1.5pt},
  flujo/.style={-{Latex[length=2mm,width=1.6mm]}, draw=fusalPrim!85, semithick},
  flujoprop/.style={-{Latex[length=2mm,width=1.6mm]}, draw=fusalSec,
    dashed, dash pattern=on 2pt off 1.5pt, semithick},
  aviso/.style={-{Latex[length=1.6mm,width=1.3mm]}, draw=fusalGris},
  etiq/.style={font=\sffamily\tiny, inner sep=1.5pt, fill=white}
}

